\section*{الفصل \ref{ch:b2:waves-on-strings} --- الأمواج على الأوتار والقضبان: معادلة دالمبير}

\begin{solution}{exo:b2:waves-on-strings:1}
$c = 2Lf = \qty{143}{m/s}$؛ و $T = \mu c^2 = \qty{61}{N}$؛ و $\lambda = 2L = \qty{1.3}{m}$ على
الوتر، و $340/110 = \qty{3.1}{m}$ في الهواء؛ وعند الدستان الخامس: $110 \times 2^{5/12} =
\qty{147}{Hz}$.
\end{solution}

\begin{solution}{exo:b2:waves-on-strings:2}
$\sqrt{E/\rho}$: الفولاذ \qty{5.1}{km/s}، والألومنيوم \qty{5.1}{km/s}، والرصاص
\qty{1.2}{km/s}. وفي \qty{1}{km}: \qty{0.20}{s} في القضيب، و \qty{2.9}{s} في الهواء ---
أي صوتان، صوت القضيب أولًا، بينهما \qty{2.7}{s}.
\end{solution}

\begin{solution}{exo:b2:waves-on-strings:3}
$400$ و $600$ و $800$ هرتز، وفيها $1$ و $2$ و $3$ عقد داخلية؛ وعقدتا النمط 3 عند
$L/3$ و $2L/3$. ومع $T \times 2$: $200\sqrt2 = \qty{283}{Hz}$؛ وعند $L/2$: \qty{400}{Hz}؛ ومع $\mu
\times 3$: $200/\sqrt3 = \qty{115}{Hz}$.
\end{solution}

\begin{solution}{exo:b2:waves-on-strings:4}
$c = \sqrt{45/0.2} = \qty{15}{m/s}$؛ و \qty{24}{m} في \qty{1.6}{s}؛ وتعود النبضة
مقلوبة ($r = -1$)؛ ومع الحلقة الحرّة تعود قائمة ($r = +1$).
\end{solution}

\begin{solution}{exo:b2:waves-on-strings:5}
(أ) $\mu = \rho\pi d^2/4 = \qty{6.1}{g/m}$؛ و $c = 2Lf = \qty{334}{m/s}$؛ و $T = \mu c^2 =
\qty{690}{N}$. (ب) $\sigma = T/A = \rho c^2 = \qty{0.87}{GPa}$، والهامش $2.3$. (ج) $T/4 =
\qty{170}{N}$، والإجهاد نفسه (إذ لا يتوقف إلا على $\rho c^2$). (د) والوتر الطويل
لا يتّسع له الصندوق؛ والسلك الصرف السميك صلب جدًّا (شديد
اللاتوافقية) ويصعب ثنيه فوق الفرس؛ ولفّ النحاس يضيف الكتلة
دون الصلابة.
\end{solution}

\begin{solution}{exo:b2:waves-on-strings:6}
(أ) $m_n = 4L^2f^2\mu/n^2g = 6.6/n^2$ كيلوغرام: أي $6.6$ و $1.65$ و $0.73$ و \qty{0.41}{kg}. (ب)
$n^2 = 13.2$، و $n = 3.6$: فلا رنين، وهي بين النمطين 3 و 4 (وأقربهما 4،
عند \qty{0.41}{kg}). (ج) عقدة عمليًّا: إذ سعة الهزّاز ضئيلة
بالقياس إلى العروات. (د) وتُغذّى الطاقة في الطور نفسه في كل دور
فتتراكم؛ ويحدّ التخامد السعةَ (الهواء، والاحتكاك الداخلي، والتسرّب
عبر البكرة).
\end{solution}

\begin{solution}{exo:b2:waves-on-strings:7}
$c = \qty{100}{m/s}$، و $k = \qty{6.3}{rad/m}$، و $\lambda = \qty{1.0}{m}$؛ و $v_{\max} = \omega A =
\qty{1.3}{m/s}$، و $a_{\max} = \omega^2A = \qty{790}{m/s^2}$؛ والميل $kA = 0.013$؛ و $Z =
\sqrt{T\mu} = \qty{1.0}{kg/s}$؛ و $\langle\mathcal P\rangle = \tfrac12Z\omega^2A^2 = \qty{0.79}{W}$؛
وفي طول الموجة الواحد $\mathcal P/f = \qty{7.9}{mJ}$.
\end{solution}

\begin{solution}{exo:b2:waves-on-strings:8}
(أ) $c_1 = \qty{100}{m/s}$، و $c_2 = \qty{50}{m/s}$؛ و $Z_1 = \qty{0.20}{kg/s}$، و $Z_2 = \qty{0.40}{kg/s}$.
(ب) من الوتر الخفيف: $r = -1/3$، و $t = 2/3$؛ ومن الثقيل: $r =
+1/3$، و $t = 4/3$. (ج) $R = 1/9$، و $\mathcal T = 4 \times 0.08/0.36 = 8/9$ في الحالتين.
(د) والنافذة: ارتفاعها \qty{0.67}{cm} وعرضها \qty{5}{cm} (المدة نفسها، ونصف
السرعة)؛ والمنعكسة: \qty{-0.33}{cm} و \qty{10}{cm}.
\end{solution}

\begin{solution}{exo:b2:waves-on-strings:9}
(أ) $E = F\ell/A\Delta\ell = 80 \times 2/(7.85 \times 10^{-7} \times 10^{-3}) = \qty{2.0e11}{Pa}$.
(ب) $c = \qty{5.1}{km/s}$. (ج) الطولي: $c/2\ell = \qty{1.3}{kHz}$؛ والمستعرض:
$\sqrt{T/\mu} = \qty{114}{m/s}$، و $f = \qty{29}{Hz}$ --- إذ يجهد التوتر
السلك حتى \qty{100}{MPa}، أي $5 \times 10^{-4}$ من $E$، ومنه تكون السرعة المستعرضة
$\sqrt{5 \times 10^{-4}}$ من الطولية. (د) $k = Ea = \qty{50}{N/m}$؛ و $m
= \rho a^3 = \qty{1.2e-25}{kg}$ (أي نحو $70$ واحدة كتلة ذرية).
\end{solution}

\begin{solution}{exo:b2:waves-on-strings:10}
(أ) $b_n = \frac2L\int_0^Ly(x, 0)\sin(n\pi x/L)\dd x$ مع المثلث $y =
2hx/L$ على $[0, L/2]$ (وبالتناظر): وبالمكاملة بالتجزئة، $b_n = 8h\sin(n\pi/2)
/n^2\pi^2$. (ب) وتنعدم من أجل $n$ الزوجي: إذ المنتصف، وهو عقدة لكل نمط زوجي،
هو النقطة المشدودة. (ج) $E_n = \tfrac14\mu L\omega_n^2b_n^2 \propto n^2/n^4 = 1/n^2$
(من أجل $n$ الفردي): وللنمط الأساسي $1/\sum_{\text{فردي}}n^{-2} = 8/\pi^2 = 81\%$. (د)
$b_n \propto \sin(n\pi/5)/n^2$: فتنعدم $n = 5, 10, \dots$. وقرب الفرس يكون
المعامل $\sin(n\pi x_0/L)$ صغيرًا في الأنماط الدنيا ويظل ينمو
في العليا: أي توافقيات عالية أكثر نسبيًّا، وصوت أنصع.
\end{solution}

\begin{solution}{exo:b2:waves-on-strings:11}
(أ) $-m\omega^2 = k(2\cos ka - 2) = -4k\sin^2(ka/2)$. (ب) $v_\varphi = 2\omega_0\sin(ka/2)
/k$، و $v_g = \omega_0a\cos(ka/2)$؛ وكلتاهما $\to \omega_0a = c$ حين $ka \to 0$؛ وعند $ka = \pi$:
$v_\varphi = 2c/\pi$، و $v_g = 0$. (ج) ويتحرك الجاران في تعاكس الطور: أي \omterm{prop:b2:waves-on-strings:modes}{موجة
مستقرة} لا تنقل شيئًا؛ و $k$ فيما وراء $\pi/a$ يعطي مجموعة الانتقالات
نفسها التي يعطيها $k' = k - 2\pi/a$ --- فلا موجة جديدة. (د) $\omega_0 = c/a =
\qty{2.0e13}{rad/s}$، و $f_{\max} = \omega_0/\pi = \qty{6.5}{THz}$؛ وعند \qty{1}{MHz}: $\lambda =
\qty{5.1}{mm} = 2 \times 10^7a$، و $ka \sim 3 \times 10^{-7}$: أي غير متبدّد إطلاقًا.
\end{solution}

\begin{solution}{exo:b2:waves-on-strings:12}
(أ) $T(x) = \mu gx$، و $c = \sqrt{gx}$. (ب) $t = \int_0^L\dd x/\sqrt{gx} = 2\sqrt{L/g} =
\qty{2.0}{s}$؛ ويمضي أعلى النبضة أسرع من أسفلها: فتتمطّ.
(ج) و $c = \sqrt{T/\mu}$ ينمو نحو الطرف. (د) الطاقة $\approx
\mathcal P\tau = ZV^2\tau = \sqrt{T\mu}V^2\tau$ ثابتة: ومنه $V \propto \mu^{-1/4}$؛ ومعامل
$10^4$ في $\mu$ يعطي $\times10$: أي \qty{100}{m/s} --- ويكسب السوط الحقيقي
الباقي من العروة المنفكّة، وهي تركّز الطاقة أكثر،
فيتجاوز الطرف سرعة الصوت.
\end{solution}

\begin{solution}{pb:b2:waves-on-strings:1}
\textbf{1.} $\mu = \rho\pi d^2/4 = \qty{6.1}{g/m}$؛ و $c = 2Lf = \qty{334}{m/s}$؛ و $T = \mu c^2
= \qty{685}{N}$.

\textbf{2.} $\sigma = \rho c^2 = \qty{0.87}{GPa}$: فالهامش $2.3$. و $T \propto f^2$: فنصف
نغمة أحدّ يضيف $12\%$، وديوان أعلى يربّعه، أي
\qty{3.5}{GPa} --- فينقطع.

\textbf{3.} $230 \times 685 \approx \qty{160}{kN}$، أي ستة عشر طنًّا.

\textbf{4.} $c = 2Lf = \qty{105}{m/s}$؛ وعند \qty{685}{N}: $\mu = T/c^2 = \qty{63}{g/m}$،
أي قضيب فولاذي قطره \qty{3.2}{mm} --- وهو أصلب من أن يهتزّ
اهتزازًا توافقيًّا أو ينثني فوق الفرس؛ ولبّ رفيع ملفوف بالنحاس
يعطي الكتلة دون الصلابة.

\textbf{5.} $440$ و $880$ و $1320$ و $1760$ و $2200$ و $2640$ و $3080$ و \qty{3520}{Hz}: فالتوافقيات $2$
و $4$ و $8$ ديوانات؛ و $3$ و $6$ خماسية فوق ديوان؛ و $5$ ثالثة كبيرة
تامة فوق ديوانين؛ و $7$ ($3080/1760 = 1.75$) سابعة صغرى
مخفوضة --- أي متنافرة.

\textbf{6.} $E = \tfrac14\mu L\omega^2A^2 = 0.25 \times 6.1 \times 10^{-3} \times 0.38 \times (2\pi
\times 440)^2 \times (5 \times 10^{-4})^2 = \qty{1.1}{mJ}$.

\textbf{7.} $A\sin kx\cos\omega t = \tfrac12A[\sin(kx - \omega t) + \sin(kx + \omega t)]$: أي موجتان
سعة كل منهما \qty{0.25}{mm}؛ و $Z = \mu c = \qty{2.0}{kg/s}$؛ وتحمل كل واحدة
$\tfrac12Z\omega^2(A/2)^2 = \qty{0.49}{W}$.

\textbf{8.} $\dot y(x, 0) = \sum_nB_n\omega_n\sin(n\pi x/L)$؛ نضرب في
$\sin(m\pi x/L)$ ونكامل، مستعملين $\int_0^L\sin\sin = \tfrac L2\delta_{mn}$.

\textbf{9.} المكمول هو $v_0\sin(n\pi x/L)$ على عرض $w$ حول $x_0$،
حيث يكون الجيب شبه ثابت.

\textbf{10.} $\sin(n\pi/8) = 0$ من أجل $n = 8, 16, \dots$؛ ولأن $B_n/B_1 =
\sin(n\pi/8)/(n\sin(\pi/8))$، ينخفض التوافقي السابع إلى $1/7$ من
الأول بينما تحتفظ التوافقيات المتآلفة من $2$ إلى $6$ بنسب من $0.9$ إلى $0.3$: فتُروَّض
السابعة المتنافرة.

\textbf{11.} تغادر نبضتا سرعة نقطةَ الضرب في جهتين متعاكستين،
وتنعكسان مقلوبتين عند الطرفين، وتتقاطعان وتتراكبان من جديد؛
وتتكرر الحالة الابتدائية بعد $2L/c = 1/f_1 = \qty{2.3}{ms}$.

\textbf{12.} $B = \pi^3 \times 2 \times 10^{11} \times 10^{-12}/(64 \times 685 \times 0.144) =
9.8 \times 10^{-4}$؛ و $\sqrt{1 + 64B} = 1.031$: أي $+3.1\%$، و $1200\log_2(1.031) = 53$ سنتًا.

\textbf{13.} توافقيات النغمة المنخفضة أحدّ من المضاعفات
التامة؛ ولكي تنطبق على الأنماط الأساسية للنغمات الأعلى
(وتُتجنَّب الضربات)، تُوزَن النغمات العالية أحدّ قليلًا.

\textbf{14.} $Z_s = \mu c = \qty{2.0}{kg/s}$؛ و $r = (2 - 2000)/2002 = -0.998$.

\textbf{15.} $1 - r^2 \approx 4Z_s/Z_b = 4 \times 10^{-3}$.

\textbf{16.} $c/2L = 440$ انعكاسًا في الثانية: فمعدل التخامد $440 \times 4 \times
10^{-3} = \qty{1.8}{s^{-1}}$، و $\tau = \qty{0.56}{s}$؛ و $\ln10^6/1.8 = \qty{7.7}{s}$ من أجل
\qty{60}{dB}.

\textbf{17.} $1.1\,\text{mJ} \times 1.8\,\text{s}^{-1} = \qty{2}{mW}$ --- أي مئتا ضعف
حديث هادئ؛ وهو ما تسمعه.

\textbf{18.} سلك قطره \qty{1}{mm} لا يكاد يدفع هواءً: إذ ينزلق الهواء حوله
وتكون استطاعته المشعّة مهمَلة. ويأخذ لوح الرنين طاقة
الوتر عبر الفرس ويحرّك الهواء على متر مربع:
أي مواءمة ممانعة بالمساحة.

\textbf{19.} إن \qty{440}{Hz} هي التوافقي الثاني للا$_3$: فيهزّ الفرس
ذلك الوتر عند أحد أنماطه --- أي رنين بالتعاطف.

\textbf{20.} ترفع $Z_b$: فتتسرّب طاقة أقلّ في كل انعكاس، وتدوم النغمة
أطول --- وتكون أخفت.

\textbf{21.} $c = \qty{5.1}{km/s}$: أي \qty{1.0}{s} في القضيب، و \qty{15}{s} في الهواء.

\textbf{22.} $50/5064 = \qty{10}{ms}$؛ وعلى الإجهاد أن ينعدم عند طرف حرّ:
فيعود الانضغاط تخلخلًا.

\textbf{23.} $\lambda = \qty{2.5}{mm}$؛ و $Z_{\text{فولاذ}} = \rho c = \qty{4e7}{kg.m^{-2}.s^{-1}}$
مقابل $400$: ومنه $r \approx -1$، أي انعكاس تام: فأي شرخ مفتوح، مهما
رقّ، يعيد صدى.

\textbf{24.} $\Delta t = d(1/3.5 - 1/6) = d \times \qty{0.119}{s/km}$: ومنه $d = \qty{170}{km}$.

\textbf{25.} الوتر \qty{334}{m/s} ($T$ و $\mu$)؛ والقضيب \qty{5.1}{km/s} ($E$ و $\rho$)؛
والانضغاط في الصخر \qty{6}{km/s} (معامل الانضغاط و $\rho$)؛ والقصّ في الصخر \qty{3.5}{km/s}
(معامل القصّ و $\rho$)؛ والهواء \qty{340}{m/s} ($\gamma P$ و $\rho$ --- في الفصل التالي).
\end{solution}
